\section{Related Work}
\paragraph{Benchmarks for Web Agents}
Early web agent benchmarks such as MiniWoB++~\citep{furuta2024multimodalwebnavigationinstructionfinetuned} are built in simplified synthetic environments. Later benchmarks, including Mind2Web~\citep{deng2023mind2web}, Mind2Web-Live~\citep{pan2024webcanvasbenchmarkingwebagents}, WebVoyager~\citep{he2024webvoyagerbuildingendtoendweb}, Online-Mind2Web~\citep{online-mind2web}, and Mind2Web~2~\citep{gou2025mind2web2evaluatingagentic}, move toward real or near-real websites, but their evaluations are still largely restricted to anonymous access without modeling authentication or session persistence. WebArena~\citep{zhou2023webarena} and VisualWebArena~\citep{koh2024visualwebarenaevaluatingmultimodalagents} support login in self-hosted environments, but are limited to small-scale sandbox websites. Moreover, in terms of task coverage, existing benchmarks are still dominated by information-retrieval tasks, with only a few considering malicious user prompts (see \cref{tab:dataset_comparison}) or transactional requirements (e.g., WebBench~\cite{webbench2025} and WorkArena~\cite{drouin2024workarenacapablewebagents}), and their task-type distributions remain highly imbalanced. Motivated by this observation, we propose LexBench-Browser, which evaluates agents on real websites with authentication and session persistence, and balances browsing-oriented needs such as search and comparison with transactional tasks, thereby systematically addressing these limitations while further extending coverage to Chinese platforms, time-varying content, and adversarial scenarios.



\paragraph{Automatic Evaluation for Web}
% Evaluation methods for web agents fall into three categories: rule-based, LLM-as-a-Judge, and hybrid approaches. Rule-based methods~\citep{} rely on predefined success conditions such as URL matching or element state verification, offering high precision but limited flexibility for tasks with multiple valid solutions. LLM-as-a-Judge methods~\citep{} leverage language models to assess task completion from screenshots or action histories, enabling more flexible evaluation but suffering from high variance across runs and substantial computational cost that scales with trajectory length. Recent work such as WebJudge~\citep{} and Agent-as-a-Judge~\citep{} attempts to improve reliability through structured prompting or multi-agent verification, yet still requires the judge to simultaneously interpret task requirements and assess execution quality. LexJudge decouples these concerns: task understanding is front-loaded to annotation time through structured rubrics, reducing evaluation to a verification task that achieves higher consistency across judge models while requiring only a single LLM call per task.

%长一些
% Unlike offline settings, online evaluation is inherently more challenging due to the involvement of real-world websites, dynamic states, and long-horizon trajectories. Early work primarily focused on enabling scalable automatic evaluation in real or near-real environments. Autonomous Evaluation~\cite{pan2024autonomousevaluationrefinementdigital} introduced MLLMs to automatically verify agent outcomes, significantly reducing human labor, but it considers only the final screenshot and ignores intermediate states, leading to substantial information loss. WebVoyager~\cite{he2024webvoyagerbuildingendtoendweb} subsequently extended evaluation to long trajectories based on full screenshot sequences, incorporating more process information in form, but its judgments remain largely holistic and suffer from severe token and computational overhead. AgentTrek~\cite{xu2025agenttrekagenttrajectorysynthesis} evaluates agent behaviors using task descriptions, actions, and reasoning traces, but still remains at the level of binary outcome judgment. As the companion evaluator of Online-Mind2Web~\cite{online-mind2web}, WebJudge mitigates information loss and token overload to some extent by identifying and retaining key intermediate screenshots, yet its evaluation signal is still predominantly holistic. Overall, existing online evaluation methods mainly rely on \emph{static criteria} and \emph{holistic judgments}, lacking explicit task decomposition, and therefore cannot provide fine-grained feedback for error analysis and capability diagnosis (see \cref{tab:evaluator_comparison}). In contrast, LexJudge follows the same online-based verification setting, but advances evaluation to an \emph{interpretable and accumulative process-level} paradigm by introducing structured scoring items and key-step alignment; meanwhile, under a fixed success criterion, it controlledly absorbs new failure patterns and valid behaviors over time, allowing the evaluator to continuously improve with usage.
Unlike offline settings, online evaluation is inherently more challenging due to real-world websites, dynamic states, and long-horizon trajectories. Existing work has mainly focused on enabling \emph{scalable automatic verification} in real or near-real environments. Representative methods differ mostly in how much trajectory information they retain: Autonomous Evaluation~\cite{pan2024autonomousevaluationrefinementdigital} verifies only the final screenshot, WebVoyager~\cite{he2024webvoyagerbuildingendtoendweb} extends this to full screenshot sequences at substantial cost, and WebJudge (the evaluator of Online-Mind2Web~\cite{online-mind2web}) keeps only key intermediate screenshots. Meanwhile, AgentTrek~\cite{xu2025agenttrekagenttrajectorysynthesis} and similar work also remain at the level of holistic binary outcome judgment. Despite these differences, existing methods still rely on \emph{static criteria} and \emph{holistic judgments}, lacking explicit task decomposition, and thus provide little diagnostic value (see \cref{tab:evaluator_comparison}). In contrast, LexJudge advances evaluation to an \emph{interpretable and accumulative process-level} paradigm via structured scoring items and key-step alignment, and can continuously improve by absorbing new failure patterns and valid behaviors over time.




% \paragraph{uni platform?}
